% 附录 D 磁学中的能量与退磁场（Energy in magnetism and demagnetizing fields）
\chapter{磁学中的能量与退磁场}

\section{能量}

计算磁化介质在磁场中的能量是件出人意料精微的事。可以得到各种表达式，但它们所指
的对象可能不同。在磁学中，这取决于你谈的是磁矩单独的能量，还是磁矩加上提供磁场
的东西二者的能量，以及这一能量究竟如何划分。

例如，把一把螺丝刀靠近大磁体的极靴，螺丝刀会被强烈吸向场最强的区域。用被磁化的
螺丝刀替换极靴间隙中的一部分空气似乎是能量上有利的事——这暗示可磁化介质（螺丝刀）
的存在降低了能量。然而，在带铁芯的磁体中建立电流比在无铁芯的磁体中要费更多能量
——这暗示可磁化介质（此处为铁芯）的存在提高了能量。那么能量究竟是降低还是升高？
答案取决于你考虑的是哪个能量。第一个例子里，我们只量度了握螺丝刀的手臂感受到的
力所省下的能量，却忽略了磁体电源为把电流维持在我们挥舞螺丝刀之前的水平而必须
多做的功。这提醒我们：处理磁性材料的能量学必须极其小心（详细处理见延伸阅读中
Heine 的文章）。

附录 B 中我们已经看到，$\delta W=-\boldsymbol{\mu}\cdot\delta\mathbf{B}$ 是磁化
介质应考虑的恰当自由能。这一表达式在统计力学中表述磁性时将有用武之地（见附录 E）。
本附录集中于：铁磁材料的存在如何因退磁场的引入而同时改变 H 与 B——退磁场自身
是要耗费能量的。

\section{退磁因子}

铁磁体内部的磁化强度 M 到达表面时必须突然停止，因此 M 有散度。利用方程
\begin{equation*}
\nabla\cdot\mathbf{H} = -\nabla\cdot\mathbf{M},\tag{D.1}
\end{equation*}
（它由 $\nabla\cdot\mathbf{B}=0$ 得来，见附录 B）我们发现 $\mathbf{H}$ 有等大反向
的散度。情形就好像铁磁体表面留下了磁单极，而这些单极充当 $\mathbf{H}$ 的源。
由此产生的 H 场称为退磁场（demagnetizing field）。图 D.1 以简单情形展示这一局面：
无限大平板或薄膜铁磁材料的磁化。若磁化位于板面内，M 的散度只出现在（设为无穷远
的）两端，因此没有退磁场。若磁化垂直于板面，顶、底表面产生磁极，板内出现退磁场
$\mathbf{H}_{\mathrm{d}}=-\mathbf{M}$。

\begin{figure}
\centering
\includegraphics[width=0.6\textwidth]{figD_01.pdf}
\caption{无限大平板的退磁（截面视图）。（a）若磁化位于板面内，板表面不产生磁极
（除两端极小的磁极外）。（b）若磁化如图垂直于板面，顶（底）面上 M 的负（正）散度
产生 H 的正（负）散度，从而如图在顶（底）面产生正（负）磁极。（c）这导致一个从
正磁极流向负磁极的退磁场。}
\end{figure}

对任意形状的铁磁体，退磁场可以是位置极其复杂的函数。但对椭球铁磁体它取相对简单
的形式：此时它在铁磁体内部是均匀的，值为
\begin{equation*}
\mathbf{H}_{\mathrm{d}} = -\mathbf{N}\mathbf{M}\tag{D.2}
\end{equation*}
N 是退磁张量。一般地可写
\begin{equation*}
(H_{\mathrm{d}})_i = -\sum_jN_{ij}M_j.\tag{D.3}
\end{equation*}
若 $\mathbf{M}$ 沿椭球的一个主轴，$\mathsf{N}_{ij}$ 可以对角化为
\begin{equation*}
\mathsf{N} = \begin{pmatrix}N_x & 0 & 0\\ 0 & N_y & 0\\ 0 & 0 & N_z\end{pmatrix}.\tag{D.4}
\end{equation*}
退磁张量满足
\begin{equation*}
\mathrm{Tr}\,\mathsf{N} = N_x+N_y+N_z = 1.\tag{D.5}
\end{equation*}

\begin{example}{D.1}
现在考察几个特例。

（1）球：$N_x=N_y=N_z=\frac{1}{3}$，故
\begin{equation*}
\mathbf{H}_{\mathrm{d}} = -\frac{\mathbf{M}}{3}\tag{D.6}
\end{equation*}

（2）沿 z 的细长圆柱棒：$N_x=N_y=\frac{1}{2}$、$N_z=0$。因为若磁化严格沿棒方向，
产生的磁极在棒的两端——可以假定远得不值得一提。

（3）垂直于 z 的平板：$N_x=N_y=0$、$N_z=1$。即上面考虑过的情形。
\end{example}

\section{任意形状的铁磁体}

本节说明如何计算外加磁场中任意形状铁磁体的退磁能。对此问题，磁场
$\mathbf{H}(\mathbf{r})$ 与 $\mathbf{B}(\mathbf{r})$ 可以分解为两个分量：
\begin{equation*}
\mathbf{H}(\mathbf{r}) = \mathbf{H}_{\mathrm{a}}(\mathbf{r}) + \mathbf{H}_{\mathrm{d}}(\mathbf{r})\tag{D.7}
\end{equation*}
\begin{equation*}
\mathbf{B}(\mathbf{r}) = \mathbf{B}_{\mathrm{a}}(\mathbf{r}) + \mathbf{B}_{\mathrm{d}}(\mathbf{r}),\tag{D.8}
\end{equation*}
外场为 $\mathbf{H}_{\mathrm{a}}(\mathbf{r})=\mathbf{B}_{\mathrm{a}}(\mathbf{r})/\mu_0$，
退磁场为 $\mathbf{H}_{\mathrm{d}}(\mathbf{r})$；把 $\mathbf{B}_{\mathrm{d}}(\mathbf{r})$
写成退磁场产生的磁通密度 $\mu_0\mathbf{H}_{\mathrm{d}}(\mathbf{r})$ 与材料内部
$\mu_0\mathbf{M}(\mathbf{r})$ 之和，即
$\mathbf{B}_{\mathrm{d}}(\mathbf{r})=\mu_0(\mathbf{H}_{\mathrm{d}}(\mathbf{r})+\mathbf{M}(\mathbf{r}))$。
退磁场源于磁体表面上凡 $\nabla\cdot\mathbf{M}(\mathbf{r})\neq0$ 处产生的磁极。

现在陈述一个重要结果。退磁场 $\mathbf{B}_{\mathrm{d}}(\mathbf{r})$ 与
$\mathbf{H}_{\mathrm{d}}(\mathbf{r})$ 满足
\begin{equation*}
\int_{\text{全空间}}\mathbf{B}_{\mathrm{d}}\cdot\mathbf{H}_{\mathrm{d}}\,\mathrm{d}\tau = 0.\tag{D.9}
\end{equation*}
证明如下：由于 $\nabla\times\mathbf{H}_{\mathrm{d}}=0$，退磁场 $\mathbf{H}_{\mathrm{d}}$
可以写成标量函数 $\phi$ 的梯度（如 $\mathbf{H}_{\mathrm{d}}=-\nabla\phi$）。于是
\begin{equation*}
\nabla\cdot(\phi\mathbf{B}_{\mathrm{d}}) \equiv \phi\nabla\cdot\mathbf{B}_{\mathrm{d}} + (\nabla\phi)\cdot\mathbf{B}_{\mathrm{d}} = (\nabla\phi)\cdot\mathbf{B}_{\mathrm{d}}\tag{D.10}
\end{equation*}
最后一步用了 $\nabla\cdot\mathbf{B}_{\mathrm{d}}=0$。从而
\begin{equation*}
\int_{\text{全空间}}\mathbf{B}_{\mathrm{d}}\cdot\mathbf{H}_{\mathrm{d}}\,\mathrm{d}\tau = -\int_{\text{全空间}}\nabla\cdot(\phi\mathbf{B}_{\mathrm{d}})\,\mathrm{d}\tau = \lim_{R\to\infty}\int_{\text{面}}\phi\,\mathbf{B}_{\mathrm{d}}\cdot\mathrm{d}\mathbf{S} = 0,\tag{D.11}
\end{equation*}
面积分取在半径 R 趋于无穷的球面上。若铁磁体尺度有限，积分必为零：$R\to\infty$ 时
至多 $B_d\sim R^{-2}$、$\phi\sim R^{-1}$。

若不存在外场，退磁场的能量 E 就是能量密度 $\frac{1}{2}\mu_0H_d^{2}$ 对全空间的
积分。利用式 D.9 可以把它化为铁磁体体积 V 内的积分，使能量表示为
\begin{align*}
E = \frac{\mu_0}{2}\int_{\text{全空间}}\mathbf{H}_{\mathrm{d}}^{2}\,\mathrm{d}\tau &= -\frac{1}{2}\int_{\text{全空间}}(\mathbf{B}_{\mathrm{d}}-\mu_0\mathbf{M})\cdot\mathbf{H}_{\mathrm{d}}\,\mathrm{d}\tau\\
&= -\frac{\mu_0}{2}\int_{\text{全空间}}\mathbf{M}\cdot\mathbf{H}_{\mathrm{d}}\,\mathrm{d}\tau\\
&= -\frac{\mu_0}{2}\int_V\mathbf{M}\cdot\mathbf{H}_{\mathrm{d}}\,\mathrm{d}\tau,\tag{D.12}
\end{align*}
最后一行来自 M=0（体外）。应当注意：尽管最后的积分只遍及体内，它表达的却是全部
杂散场（包括体外杂散场）的能量。

若存在外场，铁磁体的能量可以写成总场能量密度 $\frac{1}{2}\mu_0H^{2}$ 与外场能量
密度 $\frac{1}{2}\mu_0H_a^{2}$ 之差——这去掉了均匀外场对全空间积分给出的无穷大
能量贡献。于是能量可以有效地写为
\begin{equation*}
E = \frac{\mu_0}{2}\int_{\text{全空间}}(\mathbf{H}^{2}-\mathbf{H}_{\mathrm{a}}^{2})\,\mathrm{d}\tau.\tag{D.13}
\end{equation*}
利用 $\nabla\times\mathbf{H}_a=0$，有
\begin{equation*}
\int_{\text{全空间}}\mathbf{H}_{\mathrm{a}}\cdot\mathbf{B}_{\mathrm{d}}\,\mathrm{d}\tau = 0.\tag{D.14}
\end{equation*}
再利用下列方程：
\begin{equation*}
\mathbf{H}^{2} = \mathbf{H}_{\mathrm{a}}^{2} + 2\mathbf{H}_{\mathrm{d}}\cdot\mathbf{H}_{\mathrm{a}} + \mathbf{H}_{\mathrm{d}}^{2},\tag{D.15}
\end{equation*}
\begin{equation*}
\mathbf{H}_{\mathrm{d}}\cdot\mathbf{H}_{\mathrm{a}} = \mathbf{H}_{\mathrm{a}}\cdot\left(\frac{\mathbf{B}_{\mathrm{d}}}{\mu_0}-\mathbf{M}\right).\tag{D.16}
\end{equation*}
并再用式 D.9，能量 E 可以写成熟悉的形式
\begin{equation*}
E = -\mu_0\int_V\mathbf{M}\cdot\mathbf{H}_{\mathrm{a}}\,\mathrm{d}\tau - \frac{\mu_0}{2}\int_V\mathbf{M}\cdot\mathbf{H}_{\mathrm{d}}\,\mathrm{d}\tau,\tag{D.17}
\end{equation*}
即塞曼项与退磁项之和，每个积分只遍及铁磁体体积 V。式 D.17 的能量尽管如此仍隐含
地包括体外退磁场相联系的能量。

\begin{example}{D.2}
退磁项中二分之一因子可以这样理解：避免重复计数。塞曼项是磁化强度与外场的相互
作用；退磁项是自能——产生场的矩与自己的场相互作用。二分之一因子保证这一能量
不被计两次。

半径 a、磁化强度均匀为 M 的铁磁球是导出式 D.12 诸想法的详细示例（见图 D.2）。
退磁场为
\begin{equation*}
\mathbf{H}_{\mathrm{d}} = -\frac{\mathbf{M}}{3}\tag{D.18}
\end{equation*}
\begin{equation*}
\mathbf{B}_{\mathrm{d}} = \frac{2\mu_0\mathbf{M}}{3}.\tag{D.19}
\end{equation*}
\begin{figure}
\centering
\includegraphics[width=0.4\textwidth]{figD_02.pdf}
\caption{磁化强度均匀为 M 的铁磁球。球内区域（体积 $V_{\text{sphere}}=\frac{4}{3}\pi a^{3}$）
记为 V，球外区域记为 $\bar{V}$。}
\end{figure}
于是球内（体积 $V_{\text{sphere}}=\frac{4}{3}\pi a^{3}$ 的空间，记为 V）简单积分
给出：
\begin{equation*}
-\frac{\mu_0}{2}\int_V\mathbf{M}\cdot\mathbf{H}_{\mathrm{d}}\,\mathrm{d}\tau = \frac{1}{6}\mu_0M^{2}V_{\text{sphere}}\tag{D.20}
\end{equation*}
\begin{equation*}
\frac{\mu_0}{2}\int_V\mathbf{H}_{\mathrm{d}}^{2}\,\mathrm{d}\tau = \frac{1}{18}\mu_0M^{2}V_{\text{sphere}}\tag{D.21}
\end{equation*}
\begin{equation*}
\frac{1}{2}\int_V\mathbf{B}_{\mathrm{d}}\cdot\mathbf{H}_{\mathrm{d}}\,\mathrm{d}\tau = -\frac{1}{9}\mu_0M^{2}V_{\text{sphere}}\tag{D.22}
\end{equation*}
球外（记为 $\bar{V}$ 的空间）$B_d=\mu_0H_d$（球外 M=0），退磁场来自球的磁偶极矩，
其径向与极角分量为
\begin{equation*}
\begin{aligned}
(H_{\mathrm{d}})_r &= 2Ma^{3}\cos\theta/3r^{3}\\
(H_{\mathrm{d}})_{\theta} &= Ma^{3}\sin\theta/3r^{3},
\end{aligned}\tag{D.23}
\end{equation*}
对球外区域积分得：
\begin{equation*}
\frac{\mu_0}{2}\int_{\bar{V}}\mathbf{H}_{\mathrm{d}}^{2}\,\mathrm{d}\tau = \frac{1}{2}\int_{\bar{V}}\mathbf{B}_{\mathrm{d}}\cdot\mathbf{H}_{\mathrm{d}}\,\mathrm{d}\tau = \frac{1}{2\mu_0}\int_{\bar{V}}\mathbf{B}_{\mathrm{d}}^{2}\,\mathrm{d}\tau = \frac{1}{9}\mu_0M^{2}V_{\text{sphere}}.\tag{D.24}
\end{equation*}
对全空间积分于是给出
\begin{equation*}
-\frac{\mu_0}{2}\int_{\text{全空间}}\mathbf{M}\cdot\mathbf{H}_{\mathrm{d}}\,\mathrm{d}\tau = \frac{\mu_0}{2}\int_{\text{全空间}}\mathbf{H}_{\mathrm{d}}^{2}\,\mathrm{d}\tau = \frac{1}{6}\mu_0M^{2}V_{\text{sphere}}\tag{D.25}
\end{equation*}
与式 D.12 一致。联立式 D.22 与 D.24 得
\begin{equation*}
\frac{1}{2}\int_{\text{全空间}}\mathbf{B}_{\mathrm{d}}\cdot\mathbf{H}_{\mathrm{d}}\,\mathrm{d}\tau = 0\tag{D.26}
\end{equation*}
与式 D.9 一致。
\end{example}

\section*{延伸阅读}

\begin{itemize}[itemsep=0.2em]
\item V.~Heine, Proc.~Camb.~Phil.~Soc. \textbf{52}, 546 (1956)，对电磁场中物体的
热力学有深刻的处理。
\item 相关论述还可见 C.~J.~Adkins,《Equilibrium thermodynamics》, CUP 1983 与
J.~R.~Waldram,《The theory of thermodynamics》, CUP 1985。
\item 静磁能与退磁因子的深入论述见 A.~Aharoni,《Introduction to the theory of
ferromagnetism》, OUP 1996。
\item I.~Brevik, J.~Phys.: Condensed Matter \textbf{7}, 8065 (1995) 对计算浸入
铁磁液体中的铁磁体能量的问题给出了有益的分析。
\end{itemize}
